\chapter{Metodología}
\label{cap:metodologia}

Este capítulo fija la notación, la partición temporal, las métricas y los modelos. Todo lo que
se reporta en el Capítulo~\ref{cap:resultados} se define aquí.

\section{Planteamiento formal}
\label{sec:met_formal}

Sea $\mathcal{S}$ el conjunto de \emph{series}, donde cada serie $s \in \mathcal{S}$ es un par
(especie, \ue{}). Para cada serie y cada día $t$ se observa

\begin{itemize}
  \item $y_s(t) \in \mathbb{R}_{\geq 0}$: el volumen desembarcado en kilogramos (la variable
  objetivo), y
  \item $\mathbf{x}_s(t) \in \mathbb{R}^{p}$: el vector de covariables oceanográficas del día
  $t$ en la zona de la \ue{} (temperatura superficial del mar, su anomalía, el índice de olas
  de calor y las variables de color del océano).
\end{itemize}

El horizonte de pronóstico es $\horizonte = 90$ días. Esta elección no es estadística sino
operativa: es el plazo acordado con COBI para que una cooperativa pueda actuar sobre el
pronóstico. La convención del proyecto es, por tanto, usar el estado del océano de hoy para
predecir la captura de dentro de tres meses,
\begin{equation}
  \hat{y}_s(t+\horizonte) = f\big(\mathbf{x}_s(t),\ \mathbf{z}_s(t),\ s\big),
  \label{eq:problema}
\end{equation}
donde $\mathbf{z}_s(t)$ agrupa las variables deterministas del calendario y los rezagos de la
propia serie que son estrictamente pasados en $t$. En la implementación, la ecuación
\eqref{eq:problema} se realiza desplazando las covariables oceanográficas $\horizonte$ días
hacia adelante, de modo que la fila del día $t$ nunca contiene información posterior a
$t-\horizonte$; los detalles de construcción están en la Sección~\ref{sec:datos_features}.

\paragraph{Temporada.} La temporada de langosta va del 15 de septiembre al 15 de febrero del
año siguiente; fuera de ese rango rige la veda \citep{nom006pesc2016}. Se define el indicador $\mathrm{ens}_s(t) \in
\{0,1\}$ (\emph{in season}) y la partición del calendario en temporadas $\mathcal{T}_1, \ldots,
\mathcal{T}_K$. El agregado que le interesa a la cooperativa es la suma de temporada
$Y_s(\mathcal{T}_k) = \sum_{t \in \mathcal{T}_k} y_s(t)$.

\paragraph{Modelo específico y modelo global.} Un \emph{modelo específico} estima una función
$f_s$ por serie usando solo los datos de $s$. Un \emph{modelo global} estima una única $f$
sobre la unión $\bigcup_s \{(\mathbf{x}_s(t), y_s(t))\}$, con la identidad de la serie como
covariable categórica. La comparación entre ambos es una de las preguntas del trabajo.

\section{Partición temporal y validación}
\label{sec:met_particion}

Toda partición respeta el orden cronológico: no se usa validación cruzada aleatoria, que en
series de tiempo filtra el futuro hacia el pasado \citep{tashman2000,bergmeir2012}.

\paragraph{Cortes de prueba.} Se reportan sistemáticamente \emph{dos} cortes:
\begin{itemize}
  \item \textbf{01-julio-2020}, el corte canónico, elegido para ser comparable con la primera
  iteración del trabajo. Deja $\sim$3 temporadas de entrenamiento y la caída post-ola-de-calor
  \emph{fuera} del entrenamiento.
  \item \textbf{01-junio-2024}, que aprovecha las temporadas 2022-2026 incorporadas en este
  trabajo. Deja $\sim$7 temporadas de entrenamiento y la caída \emph{dentro} del entrenamiento.
\end{itemize}
Reportar ambos no es redundante: la diferencia entre columnas aísla el efecto de \emph{cuántas
temporadas} ve el modelo, que es precisamente la hipótesis bajo prueba. Cada experimento acepta
el corte como parámetro (variable de entorno \texttt{FF\_CUT\_DATE}), de modo que las dos
columnas provienen del mismo código.

\paragraph{Validación durante el ajuste.} Dentro del conjunto de entrenamiento se usa una
ventana expansiva: los datos se dividen en $n$ bloques cronológicos y el bloque $k$ se valida
con un modelo ajustado sobre todo lo anterior, de manera que cada partición de entrenamiento es
un superconjunto de la anterior y el bloque de validación siempre está en el futuro respecto al
de ajuste (Figura~\ref{fig:splits}).

\begin{figure}[htbp]
\centering
\includegraphics[width=0.62\linewidth]{1_qvdnPF8ETV9mFdMT0Y_BBA.png}
\caption{Esquema de validación con ventana expansiva. Cada partición de entrenamiento (azul)
es un superconjunto de la anterior y el bloque de validación (naranja) siempre está en el
futuro respecto a los datos con los que se ajusta el modelo.}
\label{fig:splits}
\end{figure}

\paragraph{Ausencia de fuga de información.} Se aplican tres reglas, verificadas con pruebas
automatizadas:
\begin{enumerate}
  \item Toda covariable oceanográfica es un desplazamiento de al menos $\horizonte$ días, ya sea
  del valor o de una media móvil que termina en $t$ y luego se desplaza.
  \item Toda estadística agregada ---medias y desviaciones de estandarización, estadísticas de
  imputación, los cuantiles del conformal--- se calcula \emph{únicamente} sobre el conjunto de
  entrenamiento (o de calibración) y se aplica al de prueba.
  \item En el ensamble, la predicción de un modelo que alimenta a otro se construye fuera de
  muestra con la ventana expansiva; usar la predicción en muestra le daría al segundo modelo
  una señal que en prueba no existe con esa calidad.
\end{enumerate}

\section{Métricas}
\label{sec:met_metricas}

Sean $y_1,\ldots,y_n$ las observaciones del conjunto de prueba y $\hat{y}_1,\ldots,\hat{y}_n$
las predicciones puntuales.

\paragraph{Error diario.}
\begin{align}
  \mathrm{MAE} &= \frac{1}{n}\sum_{i=1}^{n} |y_i - \hat{y}_i|, &
  \mathrm{RMSE} &= \sqrt{\frac{1}{n}\sum_{i=1}^{n} (y_i - \hat{y}_i)^2 }.
\end{align}
El MAE está en kilogramos y es interpretable; el RMSE penaliza más los errores grandes, que en
esta serie ocurren en los días de apertura.

\paragraph{Error porcentual simétrico.} Como la escala varía en dos órdenes de magnitud entre
\UEs{}, se agrega el sMAPE \citep{hyndman2006}, acotado en $[0, 200]$ y definido con la convención de que los
términos $0/0$ (día sin captura predicho como cero) cuentan como cero:
\begin{equation}
  \mathrm{sMAPE} = \frac{100}{n} \sum_{i=1}^{n}
  \frac{2\,|y_i - \hat{y}_i|}{|y_i| + |\hat{y}_i|}.
\end{equation}

\paragraph{Error en la suma de temporada.} Es la métrica que le importa a la cooperativa:
\begin{equation}
  e(\mathcal{T}_k) = 100 \cdot
  \frac{\sum_{t \in \mathcal{T}_k}\hat{y}(t) - \sum_{t \in \mathcal{T}_k} y(t)}
       {\sum_{t \in \mathcal{T}_k} y(t)} .
\end{equation}
Se reporta con signo, porque sobre-predecir y sub-predecir tienen consecuencias operativas
distintas.

\paragraph{Diagnóstico de forma.} En una serie donde la mayoría de los días están en veda, un
pronóstico casi plano obtiene un MAE bajo sin informar nada. Para distinguir el ajuste real del
aplanamiento se reportan dos diagnósticos: la razón de dispersión
$\mathrm{sd}(\hat{y})/\mathrm{sd}(y)$ ---cercana a $0$ delata un pronóstico plano, muy por
encima de $1$ un modelo que sobre-reacciona--- y la correlación de Pearson entre $\hat{y}$ y
$y$.

\paragraph{Métricas probabilísticas.} Para un cuantil $\tau \in (0,1)$ con predicción
$\hat{q}_\tau$, la pérdida \emph{pinball} es
\begin{equation}
  L_\tau(y, \hat{q}_\tau) =
  \begin{cases}
    \tau\,(y - \hat{q}_\tau) & \text{si } y \geq \hat{q}_\tau,\\
    (1-\tau)\,(\hat{q}_\tau - y) & \text{en otro caso.}
  \end{cases}
\end{equation}
El CRPS \citep{gneiting2007} se aproxima a partir de una rejilla de $m$ cuantiles mediante su
descomposición en pérdidas pinball,
\begin{equation}
  \mathrm{CRPS} \approx \frac{2}{m}\sum_{j=1}^{m} L_{\tau_j}(y, \hat{q}_{\tau_j}),
\end{equation}
y está en kilogramos, por lo que solo es comparable entre series de escala similar. Para un
intervalo $[\ell_i, u_i]$ de nivel nominal $1-\alpha$ se reportan la \textbf{cobertura
empírica} y el \textbf{ancho medio}:
\begin{align}
  \widehat{C} &= \frac{1}{n}\sum_{i=1}^{n} \mathbf{1}\{\ell_i \leq y_i \leq u_i\}, &
  \overline{W} &= \frac{1}{n}\sum_{i=1}^{n} (u_i - \ell_i).
\end{align}
Las dos se leen juntas: la cobertura sola se maximiza con un intervalo infinitamente ancho.

\section{Modelos}
\label{sec:met_modelos}

\subsection{Modelos de referencia}
ARIMA \citep{box2015} y Prophet \citep{taylor2017} sirven de piso: son autorregresivos y no usan covariables
oceanográficas, de modo que superarlos es la condición mínima para afirmar que las variables
del océano aportan señal. Los órdenes $(p,d,q)$ de ARIMA se eligen por AIC sobre el conjunto de
entrenamiento en una rejilla acotada.

\subsection{Ensambles de árboles}
XGBoost \citep{chen2016} y LightGBM \citep{ke2017} son las familias de referencia para datos
tabulares: manejan no linealidades, interacciones y valores faltantes de forma nativa, lo que
importa porque las covariables oceanográficas tienen huecos.

\subsection{Modelos recurrentes}
La LSTM \citep{hochreiter1997,greff2017,staudemeyer2019} procesa la serie como secuencia. Se usa una ventana
de 10 pasos y dos capas apiladas con \emph{dropout}. Al ser un modelo de la primera iteración
del trabajo, se conserva su arquitectura para poder comparar contra ella.

\subsection{Modelo global con la variable objetivo normalizada}
\label{sec:met_global}
El modelo global agrupa todas las series con la identidad de la serie codificada como variable
categórica. El problema del agrupamiento directo es que la función de pérdida queda dominada
por las series de mayor tonelaje. Para evitarlo se transforma la variable objetivo con
\begin{equation}
  \tilde{y} = \log(1+y),
\end{equation}
una transformación global, invertible y \emph{sin parámetros estimados} ---por lo que no puede
filtrar información del conjunto de prueba---, y se invierte con $\exp(\tilde{y})-1$ antes de
calcular cualquier métrica en kilogramos. Se comparó también la estandarización por serie
($z$-score con media y desviación calculadas solo en entrenamiento), que resulta inferior.

\subsection{Transformer de fusión temporal}
El TFT \citep{lim2021,vaswani2023} es un modelo de atención multi-horizonte con selección de
variables, diseñado para series con covariables estáticas, conocidas y observadas. Se incluye
como \emph{prueba de techo}: no como candidato a producto, sino para responder si la capacidad
adicional compra algo con este volumen de datos.

\section{Pronóstico probabilístico: regresión cuantílica conformalizada}
\label{sec:met_cqr}

El producto no es un número sino un intervalo con cobertura verificable. El procedimiento
combina regresión cuantílica con predicción conformal por partición (\emph{Mondrian}), y opera
en la escala logarítmica de la Sección~\ref{sec:met_global}:

\begin{enumerate}
  \item \textbf{Partición.} El conjunto anterior al corte se divide cronológicamente en
  \emph{entrenamiento} $\mathcal{D}_{\text{tr}}$ y \emph{calibración}
  $\mathcal{D}_{\text{cal}}$, disjuntos. La calibración se restringe a días en temporada,
  porque es el régimen en el que el intervalo se va a usar.

  \item \textbf{Cuantiles base.} Sobre $\mathcal{D}_{\text{tr}}$ se ajusta un modelo global
  cuantílico (\emph{gradient boosting} con pérdida pinball) para una rejilla de cuantiles
  $\tau_1 < \cdots < \tau_m$, obteniendo $\hat{q}_{\tau}(\cdot)$. Si los cuantiles estimados se
  cruzan, se reordenan de forma monótona por fila \citep{chernozhukov2010}.

  \item \textbf{Puntuaciones de no conformidad.} Para el intervalo de nivel $1-\alpha$ con
  cuantiles inferior $\tau_{\text{lo}} = \alpha/2$ y superior $\tau_{\text{hi}} = 1-\alpha/2$,
  se calcula en cada punto de calibración
  \begin{equation}
    E_i = \max\big\{\, \hat{q}_{\tau_{\text{lo}}}(\mathbf{x}_i) - \tilde{y}_i,\;
                       \tilde{y}_i - \hat{q}_{\tau_{\text{hi}}}(\mathbf{x}_i) \,\big\},
    \label{eq:cqr_score}
  \end{equation}
  que es positiva cuando la observación cae fuera del intervalo y negativa cuando cae dentro
  con holgura \citep{romano2019}.

  \item \textbf{Corrección.} Sea $Q_{1-\alpha}(E)$ el cuantil empírico $(1-\alpha)(1+1/|
  \mathcal{D}_{\text{cal}}|)$ de las puntuaciones. El intervalo conformalizado es
  \begin{equation}
    \Big[\, \hat{q}_{\tau_{\text{lo}}}(\mathbf{x}) - Q_{1-\alpha}(E),\;\;
            \hat{q}_{\tau_{\text{hi}}}(\mathbf{x}) + Q_{1-\alpha}(E) \,\Big],
  \end{equation}
  que bajo intercambiabilidad tiene cobertura marginal de al menos $1-\alpha$ en muestra
  finita. Se probó también una variante \emph{normalizada}, en la que la corrección se ensancha
  proporcionalmente a la dispersión local en lugar de ser constante.

  \item \textbf{Mondrian por serie.} Los pasos 3 y 4 se ejecutan \emph{por serie}: cada
  especie $\times$ \ue{} tiene su propio $Q_{1-\alpha}(E)$. Esto es indispensable porque las
  escalas difieren en órdenes de magnitud, y además permite auditar la cobertura serie por
  serie ---la cobertura marginal promedia y puede ocultar una serie mal calibrada, como se verá
  en la Sección~\ref{sec:res_cqr}.

  \item \textbf{Regreso a kilogramos.} Los extremos se invierten con $\exp(\cdot)-1$ y se
  truncan en cero. La inversión es monótona, de modo que preserva la cobertura.
\end{enumerate}

\section{Selección de hiperparámetros}
\label{sec:met_hiperparametros}

Para los modelos estadísticos se usa búsqueda exhaustiva en una rejilla acotada con criterio
AIC. Para los modelos de aprendizaje de máquina y profundo se usa optimización bayesiana con
Optuna \citep{akiba2019}, con un número limitado de ensayos para que el experimento sea
reproducible en un tiempo razonable, y evaluando siempre sobre la ventana expansiva de
validación descrita en la Sección~\ref{sec:met_particion}, nunca sobre el conjunto de prueba.
Toda componente estocástica ---inicialización de pesos, submuestreo, semillas del optimizador---
usa una semilla fija que se registra junto con las métricas del experimento.

\section{Reproducibilidad}
\label{sec:met_reproducibilidad}

El trabajo se implementó como un paquete de Python ---sobre \texttt{numpy} \citep{harris2020}, \texttt{pandas} \citep{mckinney2010}, \texttt{scikit-learn} \citep{pedregosa2011} y \texttt{PyTorch} \citep{paszke2019}--- con pruebas automatizadas
(\texttt{137} pruebas unitarias que cubren el proceso de datos, las métricas, la construcción
de variables sin fuga de información y la envoltura conformal). Cada experimento:

\begin{itemize}
  \item se ejecuta con \emph{un} comando y acepta el corte de prueba como parámetro;
  \item escribe sus métricas en un archivo JSON versionado por corte, sus figuras en PNG y un
  resumen en Markdown, de modo que ninguna cifra de este documento se transcribe a mano;
  \item fija y registra la semilla aleatoria.
\end{itemize}

Las decisiones de diseño no triviales ---el tratamiento de la variable objetivo faltante y la
justificación de la prueba de techo con el TFT--- están registradas como ADR (\emph{Architecture
Decision Records}) en el repositorio, escritas \emph{antes} de correr los experimentos
correspondientes.
